\documentclass{article}
\usepackage[utf8]{inputenc}
\usepackage{hyperref}
\usepackage{booktabs}
\usepackage{longtable}
\usepackage{array}
\usepackage{fancyvrb}
\usepackage{xcolor}
\usepackage{listings}
\usepackage{geometry}
\usepackage{graphicx}

\geometry{margin=1in}

% Code highlighting styling
\lstset{
  basicstyle=\ttfamily\small,
  breaklines=true,
  backgroundcolor=\color{gray!10},
  frame=single,
  rulecolor=\color{gray!30},
  xleftmargin=10pt,
  xrightmargin=10pt,
  showstringspaces=false
}

\title{\textbf{ReGrade --- Comprehensive LMS Regrading Request System} \\ Technical Architecture \& Specifications}
\author{UPES Development Team}
\date{June 4, 2026}

\begin{document}

\maketitle

\tableofcontents
\newpage

\section{Executive Summary \& Project Goals}

ReGrade is a lightweight, pure client-side Single Page Application (SPA) designed to handle grade disputes systematically. In academic institutions like \textbf{UPES (University of Tomorrow)}, grade disputes are often handled over unstructured, non-auditable channels like emails or direct messages. This leads to missing context, lack of accountability, and grading stress for both students and instructors.

ReGrade addresses this by replacing chaotic communications with a structured, state-governed, and auditable regrading workflow.

\subsection{Architectural Philosophy}
ReGrade is built entirely as a **client-side Single Page Application**. 
\begin{itemize}
    \item \textbf{Zero Server Dependency}: The application runs directly inside any modern web browser.
    \item \textbf{Offline Capability}: It can be executed directly from the local file system (using the \texttt{file://} protocol) without requiring CORS-configured server execution.
    \item \textbf{Local Persistence}: State is stored inside browser \texttt{localStorage} to simulate a persistent database, and session state resides in \texttt{sessionStorage}.
\end{itemize}

\begin{center}\rule{0.5\linewidth}{0.5pt}\end{center}

\section{Technology Stack \& Technical Rationale}

\subsection{Frontend Architecture}
The system is composed of standard web technologies without external frameworks or compilation steps:

\subsubsection{1. HTML5 (App Shell)}
The file \texttt{index.html} acts as the core application shell. It imports stylesheets and scripts, and sets up the structural layout:
\begin{itemize}
    \item \textbf{Sidebar Nav Container}: Renders a responsive sidebar displaying logo, role-based navigation tabs, and active user profile details.
    \item \textbf{Main Page Target}: A single container element (\texttt{<main id="page-content">}) where Javascript views inject dynamic HTML based on active routes.
    \item \textbf{Notification Drawer}: A hidden overlay panel containing the feed of read/unread notifications, sliding in from the right when triggered.
    \item \textbf{Toast Alerts}: An overlay stack container positioned at the top-right corner to catch and display status updates in real-time.
\end{itemize}

\subsubsection{2. CSS3 (Design System)}
Styling is structured across three core files inside the \texttt{css/} directory:
\begin{itemize}
    \item \textbf{\texttt{css/index.css}}: Establishes the institutional visual tokens. It defines CSS Custom Properties (variables) for HSL-based colors (deep space blue palette for dark mode), typography (Inter font family), shadows, and viewport resets.
    \item \textbf{\texttt{css/components.css}}: Houses styles for reusable elements like cards, primary/success/danger buttons, status badges, forms, timeline nodes, and the notification drawer.
    \item \textbf{\texttt{css/views.css}}: Sets up the layout structure (flex grids) for the dashboard, queue listing, login screens, and the regrading request forms.
\end{itemize}

\subsubsection{3. JavaScript (ECMAScript 6 Vanilla)}
All application routing, database management, session authentication, and view rendering are written in pure JavaScript. By design, the project avoids ES Modules (\texttt{import}/\texttt{export}) in favor of modular IIFE namespaces bound to the global \texttt{window} object. This layout allows the app to be loaded directly via \texttt{file://} in browsers that block local module imports.

\subsection{Backend \& Database Simulation}
\begin{important}
\textbf{No Java, Node, or Python backend is used for execution logic.} The entire server layer is simulated inside the client's browser.
\end{important}

\begin{itemize}
    \item \textbf{Mock Database Layer} (\texttt{js/store.js}): Simulates database tables using key-value storage in \texttt{localStorage}. It handles queries, insertions, status updates, and immutable audit logs.
    \item \textbf{Session Simulation} (\texttt{js/auth.js}): Simulates login state and role-based permissions using \texttt{sessionStorage}.
    \item \textbf{HTTP Hosting}: Any basic HTTP daemon (e.g. Python's \texttt{http.server} or Node's \texttt{npx http-server}) can serve the application.
\end{itemize}

\begin{center}\rule{0.5\linewidth}{0.5pt}\end{center}

\section{Detailed Directory \& File Map}

\begin{longtable}{llp{7cm}}
\caption{System Directory Map} \\
\toprule
\textbf{Path} & \textbf{Type} & \textbf{Technical Responsibility} \\
\midrule
\endhead
\texttt{index.html} & File & Core HTML5 application shell & layout structures. \\
\texttt{README.md} & File & Developer setup, schema guide, and walkthrough documentation. \\
\texttt{css/index.css} & File & CSS tokens, dark mode variables, global reset, utility classes. \\
\texttt{css/components.css} & File & Stylings for buttons, badges, notifications, drawer, and input fields. \\
\texttt{css/views.css} & File & Viewport grid systems, login views, queue cards, and details layouts. \\
\texttt{js/app.js} & File & SPA orchestrator, hash-based router, navigation, and notification binding. \\
\texttt{js/store.js} & File & Database abstraction layer, localStorage wrapper, database CRUD, and seed data. \\
\texttt{js/auth.js} & File & Active session simulator, role login/logout handling, and route guards. \\
\texttt{js/audit.js} & File & Immutable append-only audit trail logger helper. \\
\texttt{js/notifications.js} & File & Notification system utility (in-app drawer rendering and toast stack). \\
\texttt{js/views/login.js} & File & Renders role picker and user selector layout. \\
\texttt{js/views/dashboard.js} & File & Renders statistical dashboard metrics for students and instructors. \\
\texttt{js/views/grade-view.js} & File & Renders student grade list with "Request Regrade" indicators. \\
\texttt{js/views/request-form.js} & File & Renders new request forms, validations, student detail screens, and reopens. \\
\texttt{js/views/queue.js} & File & Renders searchable, filterable queue cards for instructors. \\
\texttt{js/views/review.js} & File & Renders instructor evaluation panels, decision entry forms, and closure actions. \\
\bottomrule
\end{longtable}

\begin{center}\rule{0.5\linewidth}{0.5pt}\end{center}

\section{State Machine \& Role Security}

\subsection{Request States}
The lifecycle of a regrade request is strictly managed in the database column \texttt{status} and follows 5 primary values:
\begin{enumerate}
    \item \textbf{\texttt{Submitted}}: A student has submitted a regrade request. The instructor can see it in their queue.
    \item \textbf{\texttt{Under Review}}: Triggered automatically when an instructor opens the request for the first time.
    \item \textbf{\texttt{Accepted}}: Grade revised upward by the instructor. The new score is written back to the student's grade record.
    \item \textbf{\texttt{Rejected}}: Original grade stands.
    \item \textbf{\texttt{Closed}}: The regrading workflow is finalized, preventing any future actions.
\end{enumerate}

\subsection{The Reopened UI Flag}
If an instructor rejects a request, the student has a **one-time dispute (reopen) guard**. 
Clicking \textbf{Reopen Request}:
\begin{itemize}
    \item Increments \texttt{reopenCount} by 1.
    \item Resets the status back to \textbf{\texttt{Submitted}}.
    \item Places a permanent `"Reopened"` flag in the UI.
    \item Blocks further reopens once \texttt{reopenCount >= 1}.
\end{itemize}

\subsection{Role-Based Navigation Guards}
Access controls are implemented client-side inside \texttt{js/app.js} and \texttt{js/auth.js}. When routing occurs, \texttt{navigate(path)} performs the following checks:
\begin{itemize}
    \item Checks if the user is authenticated via \texttt{LMS.Auth.isAuthenticated()}.
    \item Restricts student access to instructor-only pages (\texttt{\#/queue}, \texttt{\#/queue/:id}, \texttt{\#/audit}) by redirecting unauthorized attempts back to the student's dashboard.
    \item Shows/hides navigation tabs in the sidebar depending on the active role.
\end{itemize}

\begin{center}\rule{0.5\linewidth}{0.5pt}\end{center}

\section{Database Schema \& API Contract}

ReGrade uses 7 distinct localStorage tables to hold application data:

\subsection{1. \texttt{lms\_users}}
Stores identity, role representation, and avatar labels:
\begin{lstlisting}[language=json]
[
  { "id": "u1", "name": "Anirudh Bhatt", "role": "student", "email": "anirudh@upes.edu", "avatar": "AB" },
  { "id": "u2", "name": "Himanshu Kadyan", "role": "student", "email": "himanshu@upes.edu", "avatar": "HK" },
  { "id": "u3", "name": "Yashasvi Dutt Sharma", "role": "student", "email": "yashasvi@upes.edu", "avatar": "YS" },
  { "id": "u4", "name": "Raygun Jose", "role": "instructor", "email": "raygun.jose@upes.edu", "avatar": "RJ" }
]
\end{lstlisting}

\subsection{2. \texttt{lms\_assignments}}
Stores assignment titles, course labels, and maximum grade score limits:
\begin{lstlisting}[language=json]
[
  {
    "id": "a1",
    "title": "Midterm Exam",
    "description": "Chapters 1--6 comprehensive assessment",
    "maxScore": 100,
    "course": "CS301 --- Data Structures",
    "publishedAt": "2026-05-25T18:00:00.000Z"
  }
]
\end{lstlisting}

\subsection{3. \texttt{lms\_grades}}
Tracks individual student grade records:
\begin{lstlisting}[language=json]
[
  { "studentId": "u1", "assignmentId": "a1", "score": 72, "maxScore": 100 }
]
\end{lstlisting}

\subsection{4. \texttt{lms\_regrade\_requests}}
Tracks active regrading requests:
\begin{lstlisting}[language=json]
{
  "id": "rr_1780520849373_zgp3y",
  "studentId": "u1",
  "assignmentId": "a1",
  "originalGrade": 72,
  "claimedGrade": 85,
  "reason": "calculation_error",
  "reasonDetail": "Question 3 partial credit was not awarded. I showed all work on page 4.",
  "evidenceType": "text",
  "evidence": "Rubric mentions 5 points for the AVL tree AVL rotations.",
  "status": "Submitted",
  "submittedAt": "2026-06-03T18:00:00.000Z",
  "updatedAt": "2026-06-03T18:00:00.000Z",
  "reopenCount": 0
}
\end{lstlisting}

\subsection{5. \texttt{lms\_decisions}}
Records decisions made by instructors:
\begin{lstlisting}[language=json]
{
  "id": "dec_67890",
  "requestId": "rr_1780520849373_zgp3y",
  "instructorId": "u4",
  "decision": "Accepted",
  "justification": "I reviewed the work on page 4. Full credit has been awarded.",
  "revisedGrade": 85,
  "decidedAt": "2026-06-03T19:00:00.000Z"
}
\end{lstlisting}

\subsection{6. \texttt{lms\_audit\_log}}
Immutable append-only record of state-changing transactions:
\begin{lstlisting}[language=json]
{
  "id": "aud_12345",
  "requestId": "rr_1780520849373_zgp3y",
  "actorId": "u1",
  "actorRole": "student",
  "action": "SUBMITTED",
  "fromState": null,
  "toState": "Submitted",
  "metadata": {},
  "timestamp": "2026-06-03T18:00:00.000Z"
}
\end{lstlisting}

\subsection{7. \texttt{lms\_notifications}}
Tracks messages sent to users' drawers:
\begin{lstlisting}[language=json]
{
  "id": "n_12345",
  "userId": "u1",
  "type": "success",
  "message": "Your regrade request for Midterm Exam has been accepted!",
  "requestId": "rr_1780520849373_zgp3y",
  "read": false,
  "createdAt": "2026-06-03T19:00:00.000Z"
}
\end{lstlisting}

\begin{center}\rule{0.5\linewidth}{0.5pt}\end{center}

\section{User Interface & Premium Design System}

\subsection{CSS Custom Property Design Tokens}
ReGrade utilizes a hybrid visual theme combining institutional branding with high-readability content surfaces:
\begin{itemize}
    \item \textbf{UPES Background Gradient}: The global canvas background, sidebar, topbar, and login shell use a 135-degree linear gradient transitioning from deep navy-indigo through purple to plum (\texttt{--bg-gradient: linear-gradient(135deg, \#0e1e82 0\%, \#571175 60\%, \#850f5f 100\%)}).
    \item \textbf{Light Workspace Container}: Upon login, the active workspace (\texttt{.app-workspace}) floats above the gradient background as an elevated white-background sheet container (\texttt{\#ffffff}) with black/slate text (\texttt{\#1e293b}) and soft, premium shadows (\texttt{rgba(0, 0, 0, 0.22)}) for a clean, modern aesthetic.
    \item \textbf{High-Contrast Status Indicators}: Status badges inside the light workspace dynamically re-map to high-contrast colors (e.g., dark amber for Submitted, dark blue for Under Review, dark green for Accepted, and dark red for Rejected) against soft light-tinted backgrounds for optimal contrast.
    \item \textbf{Primary Action Buttons}: In order to maintain absolute visibility and focus, core interactive buttons (such as \texttt{.btn-primary}, \texttt{.btn-success}, and \texttt{.btn-danger}) are explicitly configured with high-contrast white text (\texttt{\#ffffff}) on their vivid colored backgrounds, overriding any general light-theme text styles.
\end{itemize}

\subsection{UI Logo Integration}
To integrate the landscape \textbf{UPES Logo} (which features dark text and tagline) onto the dark blue shell background:
\begin{itemize}
    \item \textbf{White Contrast Shield}: Logo containers in both the sidebar (\texttt{index.html}) and the login card (\texttt{js/views/login.js}) wrap the image inside a clean, solid white container:
    \begin{lstlisting}[language=html]
<div style="background: #ffffff; padding: 8px 16px; border-radius: 6px; box-shadow: 0 2px 8px rgba(0,0,0,0.15);">
  <img src="upes_logo.png" alt="UPES Logo" style="max-height: 38px; max-width: 100%; object-fit: contain;">
</div>
    \end{lstlisting}
    This protects the logo's legibility, ensuring the black text of the brand remains readable.
\end{itemize}

\begin{center}\rule{0.5\linewidth}{0.5pt}\end{center}

\section{Verification, Testing \& Walkthrough Script}

\subsection{Application Execution}
Since there is no complex backend compiling code, the project can be served using any local web hosting utility.

\subsubsection{Method A: Python 3 server}
\begin{lstlisting}[language=bash]
python3 -m http.server 3456
\end{lstlisting}

\subsubsection{Method B: Node.js http-server}
\begin{lstlisting}[language=bash]
npx http-server -p 3456
\end{lstlisting}

\subsection{Automated Unit Tests}
To satisfy the Day 4 testing requirements, a dedicated client-side test suite is integrated:
\begin{itemize}
    \item \textbf{File}: \texttt{test.html}
    \item \textbf{Framework}: A custom, lightweight, browser-based unit testing framework with isolated state setup.
    \item \textbf{Test Coverage}:
    \begin{itemize}
        \item \textbf{Database \& Mock Store}: Validates seed status and user configurations (including UPES emails).
        \item \textbf{Authentication \& Routing}: Verifies session initialization, logging in, and role-based guards.
        \item \textbf{Request Validation}: Confirms that valid requests are correctly registered and that duplicate request guards block multiple submissions for the same assignment.
        \item \textbf{Instructor Evaluation}: Tests status progression (Submitted $\to$ Under Review), decision logging, and gradebook updates.
        \item \textbf{Reopen Limits}: Enforces the strict single-reopen constraint (allowing exactly one reopen for rejected requests and blocking any subsequent attempts).
    \end{itemize}
    \item \textbf{Execution}: Navigate to \url{http://localhost:3456/test.html} to view the live execution metrics.
\end{itemize}

\subsection{Walkthrough Evaluation Steps}
Perform the following sequence to verify all features, states, and user names:

\subsubsection{Step 1: Student Login}
\begin{enumerate}
    \item Open your browser and navigate to \url{http://localhost:3456}.
    \item You will see the login page featuring the white-backed \textbf{UPES Logo}.
    \item Click \textbf{Student} and choose the account **Anirudh Bhatt**.
    \item Verify that the student's dashboard loads successfully, displaying their course metrics.
\end{enumerate}

\subsubsection{Step 2: Submit a New Request}
\begin{enumerate}
    \item In the sidebar, click \textbf{My Grades}.
    \item Locate **Final Project** (which currently has no request filed).
    \item Click \textbf{Request Regrade}.
    \item On the request form page:
    \begin{itemize}
        \item Claimed grade: Enter a value larger than original (e.g. 125).
        \item Reason: Select "Calculation / Addition Error".
        \item Explanation: Enter at least 30 characters (e.g. "I scored 118 out of 150 but Question 4 score was not added in my total sheet.").
        \item Evidence: Paste text details or toggle to link input.
    \end{itemize}
    \item Click \textbf{Submit Request}. Observe the confirmation toast and verification timeline.
\end{enumerate}

\subsubsection{Step 3: Instructor Evaluation}
\begin{enumerate}
    \item Click \textbf{Switch Role} at the bottom left.
    \item Select \textbf{Instructor} and select **Raygun Jose**.
    \item Navigate to the \textbf{Request Queue} in the sidebar.
    \item Observe the newly submitted request from **Anirudh Bhatt**.
    \item Click the card. The status instantly transitions from \textbf{\texttt{Submitted}} to \textbf{\texttt{Under Review}} in the background, logging an entry to the audit trail.
    \item In the decision panel: select **Accept**, enter revised grade, write a 20+ character justification, and click submit.
    \item Switch back to student **Anirudh Bhatt** to verify their updated grade, success toast, and timeline notifications.
\end{enumerate}

\end{document}
